% [EN] Build: cd docs/reports/reconstruction && latexmk -pdf synthetic_reco_B115T3deg.tex / [CN] 编译: cd docs/reports/reconstruction && latexmk -pdf synthetic_reco_B115T3deg.tex
\documentclass[aspectratio=169]{beamer}
\usetheme{Madrid}
\usecolortheme{default}

\usepackage{graphicx}
\usepackage{amsmath}
\usepackage{booktabs}
\usepackage{siunitx}

\title[Synthetic p/n Reconstruction]{Synthetic Proton/Neutron Reconstruction Study}
\subtitle{B = 1.15 T, Beam Deflection = 3.0$^\circ$}
\author{SMSimulator5.5 Analysis}
\date{\today}

\begin{document}

\begin{frame}
\titlepage
\end{frame}

\begin{frame}{Study Goal and Configuration}
\textbf{Goal}: Evaluate reconstruction quality with synthetic beam input and focus on
\[
\Delta P_x = P_x^{\mathrm{reco}} - P_x^{\mathrm{true}} .
\]

\vspace{0.3cm}
\textbf{Configuration}
\begin{itemize}
\item Magnetic setting from macro: \texttt{run\_protons\_B115T\_3.0deg.mac}
\item Geometry/field macro: \texttt{geometry\_B115T.mac}
\item Target from macro: \((-12.488, 0.009, -1069.458)\,\mathrm{mm}\), angle \(3.0^\circ\)
\item Dipole field map: \texttt{180703-1,15T-3000.table}, rotation \(30^\circ\)
\end{itemize}
\end{frame}

\begin{frame}{Synthetic Input Dataset}
\textbf{Generator setup}
\begin{itemize}
\item Proton momentum: \((x_p, 0, 627)\,\mathrm{MeV}/c\), \(x_p \sim \mathcal{U}[-150,150]\)
\item Neutron momentum: \((x_n, 0, 627)\,\mathrm{MeV}/c\), \(x_n \sim \mathcal{U}[-150,150]\)
\item Generated events: 20,000
\item Reconstruction analysis sample: 3,000 events
\end{itemize}

\vspace{0.2cm}
\textbf{Event statistics in analyzed sample}
\begin{itemize}
\item Events with reconstructed PDC track: 2597 / 3000 (\(\sim 86.6\%\))
\item Events with neutron TOF reconstruction: 994 / 3000 (\(\sim 33.1\%\))
\item Proton Minuit attempts: 600 events, Minuit success: 344 (\(\sim 57.3\%\))
\end{itemize}
\end{frame}

\begin{frame}{Reconstruction Pipelines}
\textbf{Proton (charged, PDC-based)}
\begin{itemize}
\item Build PDC track from PDC1/PDC2 reconstructed points.
\item Method A: \textbf{TMinuit} back-propagation in field to target.
\item Method B: \textbf{Three-point constrained} fit (target + two PDC points).
\end{itemize}

\textbf{Neutron (neutral, NEBULA TOF)}
\begin{itemize}
\item Cluster NEBULA hits in time window.
\item Reconstruct hit position and TOF.
\item Start point fixed at target (from macro), direction \(=\) target \(\rightarrow\) hit.
\item Compute \(\beta\), \(\gamma\), and 3D momentum:
\[
\beta=\frac{L/t}{c},\quad p=\gamma m_n \beta c .
\]
\end{itemize}
\end{frame}

\begin{frame}{PDC Position Precision}
\begin{center}
\includegraphics[width=0.92\textwidth]{../../../results/synthetic_reco/B115T_3deg/plots/pdc_position_residuals.png}
\end{center}

\small
In this synthetic setup, PDC residual histograms peak at 0 mm in the analyzed events
(\( \langle \Delta r_{PDC1} \rangle = 0 \), \( \langle \Delta r_{PDC2} \rangle = 0 \)).
\end{frame}

\begin{frame}{Proton \( \Delta P_x \): Key Result}
\begin{columns}[T]
\column{0.50\textwidth}
\includegraphics[width=\textwidth]{../../../results/synthetic_reco/B115T_3deg/plots/proton_dpx_comparison.png}

\column{0.50\textwidth}
\includegraphics[width=\textwidth]{../../../results/synthetic_reco/B115T_3deg/plots/proton_dpx_vs_true.png}
\end{columns}

\vspace{0.25cm}
\small
\begin{tabular}{lccc}
\toprule
Method & Count & Mean(\(\Delta P_x\)) & RMS(\(\Delta P_x\)) \\
\midrule
TMinuit & 600 & +230.969 MeV/c & 54.115 MeV/c \\
Three-point & 600 & 0.000 MeV/c & 0.000 MeV/c \\
\bottomrule
\end{tabular}
\end{frame}

\begin{frame}{Proton \( P_x^{true} \) vs \( P_x^{reco} \)}
\begin{center}
\includegraphics[width=0.92\textwidth]{../../../results/synthetic_reco/B115T_3deg/plots/proton_px_true_vs_reco.png}
\end{center}

\small
The Minuit branch shows a clear positive \(P_x\) bias in this sample.
\end{frame}

\begin{frame}{Neutron \( \Delta P_x \) and 3D Momentum}
\begin{columns}[T]
\column{0.48\textwidth}
\includegraphics[width=\textwidth]{../../../results/synthetic_reco/B115T_3deg/plots/neutron_dpx.png}

\column{0.52\textwidth}
\includegraphics[width=\textwidth]{../../../results/synthetic_reco/B115T_3deg/plots/neutron_px_true_vs_reco_and_dpx_vs_true.png}
\end{columns}

\vspace{0.2cm}
\small
Neutron TOF reconstruction summary (994 events):
\[
\text{Mean}(\Delta P_x)=+1.758~\mathrm{MeV}/c,\quad
\text{RMS}(\Delta P_x)=17.152~\mathrm{MeV}/c .
\]
\end{frame}

\begin{frame}{Conclusions and Next Actions}
\textbf{Conclusions}
\begin{itemize}
\item The pipeline reconstructs proton and neutron 3D momentum from synthetic input.
\item Neutron TOF reconstruction gives moderate \(P_x\) bias and finite spread.
\item Proton Minuit currently shows large positive \( \Delta P_x \) bias in this sample.
\end{itemize}

\textbf{Immediate next actions}
\begin{itemize}
\item Diagnose proton Minuit bias with event-level \((P_x^{true}, P_x^{reco})\) outliers.
\item Re-check coordinate/sign conventions between beam truth and reconstruction frame.
\item Add per-bin \(P_x^{true}\) calibration/validation before production usage.
\end{itemize}
\end{frame}

\end{document}
